\addcontentsline{toc}{subsubsection}{Solution 25. Vietnam Mathematical Olympiad 1991}

\begin{dlsummary}[title={\bfseries Solution 25. Vietnam Mathematical Olympiad 1991},breakable]

\begin{proof}
\textbf{Step 1. Determine $f(0)$.}

Setting $x=y=z=0$ gives
\[
f(0)-f(0)^2\geq\frac14,
\]
or equivalently,
\[
\left(f(0)-\frac12\right)^2\leq0.
\]
Hence
\[
\boxed{f(0)=\frac12.} \tag{1}
\]

\medskip
\noindent
\textbf{Step 2. Establish a key inequality.}

Setting $z=0$ and using (1), we obtain
\[
\frac{f(xy)}2+\frac14-\frac{f(x)}2
\geq\frac14.
\]
Therefore,
\[
\boxed{f(xy)\geq f(x)}
\qquad(x,y\in\mathbb R). \tag{2}
\]

\medskip
\noindent
\textbf{Step 3. Prove that $f$ is constant.}

For any $t\neq0$, taking $x=1,y=t$ in (2)
gives
\[
f(t)\geq f(1).
\]
Taking $x=t,y=1/t$ gives
\[
f(1)\geq f(t).
\]
Thus
\[
f(t)=f(1)\qquad(t\neq0). \tag{3}
\]

Finally, setting $x=y=z=1$ gives
\[
f(1)-f(1)^2\geq\frac14,
\]
so
\[
\left(f(1)-\frac12\right)^2\leq0.
\]
Consequently, $f(1)=\frac12$.

Combining this with (1) and (3), we get
\[
\boxed{f(x)=\frac12\qquad(x\in\mathbb R).}
\]

Conversely, for $f(x)=\frac12$, we have
\[
\frac{f(xy)+f(xz)}2-f(x)f(yz)
=\frac12-\frac14=\frac14.
\]
Thus the inequality holds with equality
for all real $x,y,z$.
\end{proof}

\end{dlsummary}